\documentclass[conference]{IEEEtran}
\IEEEoverridecommandlockouts
\usepackage{cite}
\usepackage{amsmath,amssymb,amsfonts}
\usepackage{algorithmic}
\usepackage{graphicx}
\usepackage{textcomp}
\usepackage{xcolor}
\usepackage{booktabs}
\usepackage{url}
\usepackage[hidelinks]{hyperref} % must be loaded last (after cite/url) for IEEEtran

% Draft placeholder for numbers not yet produced (first run pending on the GPU box).
\newcommand{\TODO}[1]{\textcolor{red}{[\textbf{TODO:} #1]}}
\newcommand{\NUM}[1]{\textcolor{blue}{#1}} % blue = placeholder value to be filled from the run

\def\BibTeX{{\rm B\kern-.05em{\sc i\kern-.025em b}\kern-.08em
    T\kern-.1667em\lower.7ex\hbox{E}\kern-.125emX}}
\begin{document}

\title{Probing and Steering for Safety-Relevant Behaviours 
in Large Language Models}

\author{\IEEEauthorblockN{Nimun Kaur Bajwa}
\IEEEauthorblockA{\textit{MSc. Artificial Intelligence}, Queen Mary University of London\\
{n.k.bajwa@se25.qmul.ac.uk}}}

\maketitle

\begin{abstract}
Activation probes, linear readouts of a language model's internal
representations, are used both to \emph{detect} safety-relevant behaviours and,
by adding their direction back into the residual stream, to \emph{steer} them.
These are distinct demands on a direction, and prior work has met them
separately and on different properties: Marks and Tegmark find that
difference-of-means directions for truth are the most \emph{causal};
Goldowsky-Dill et al. find that logistic-regression probes are the best
\emph{detectors} of strategic deception; McKenzie et al. find that
sequence-pooled probes detect high-stakes interactions on par with far larger
monitors. Left open is whether, for a given safety behaviour, the direction that
best detects it is also the one that best steers it, or whether there is a
\emph{read/write dissociation} in which logistic-regression directions detect
best while mass-mean directions steer best. This thesis investigates that
question across a range of safety behaviours in Llama-3.1-8B-Instruct. To keep
the comparison uniform across behaviours, each behaviour is described by one
specification (its data, token masking, readout format and steering mode) and
run through a single shared pipeline: layer selection by complementary criteria,
a grid of probe architectures, in- and out-of-distribution evaluation, and
calibrated activation steering scored by both a normalized indirect effect and
an LLM judge that separately rates behaviour and coherence. This draft reports
the results on two behaviours, deception and high-stakes awareness; the thesis
then extends the same study to sycophancy, user intention and refusal, and asks
which sparse-autoencoder latents carry each behaviour.
\TODO{tighten once results are in; state the headline finding in one sentence.}
\end{abstract}

\begin{IEEEkeywords}
mechanistic interpretability, linear probes, activation steering, AI safety,
deception detection, high-stakes monitoring, sparse autoencoders
\end{IEEEkeywords}

\section{Introduction}

As language models are deployed in consequential settings, cheap and reliable
monitoring of their internal state becomes a practical safety requirement.
A promising family of monitors are \emph{activation probes}: linear readouts
of the residual stream that flag a property of interest at a fraction of the
cost of a full auxiliary model~\cite{alain2018,mckenzie2025}. Two questions
determine whether such a probe is useful. First, does it \emph{detect} the
property; that is, does its score separate positive from negative examples,
in-distribution and, crucially, out-of-distribution? Second, is the direction
it identifies \emph{causally implicated} in the behaviour; that is, does adding
it to the residual stream during generation actually change what the model does?
These are distinct properties, and the interpretability literature has largely
studied them apart.

Three lines of work anchor this thesis. Marks and Tegmark~\cite{marks2024},
studying truth, find that simple difference-of-means (``mass-mean'') directions
are the \emph{most causal} (best at flipping the model's treatment of a
statement), even though logistic regression classifies at least as accurately.
Goldowsky-Dill et al.~\cite{goldowsky2025} (henceforth Apollo), studying
strategic deception, find that logistic-regression probes are the strongest
\emph{detectors} across diverse datasets. McKenzie et al.~\cite{mckenzie2025}
(henceforth the high-stakes work), studying whether an interaction ``might lead
to significant harm,'' find that probes which \emph{pool over the whole
sequence} (mean or learned attention) far outperform last-token probes, and
rival prompted or finetuned monitors many orders of magnitude more expensive.

Read together, these results suggest, but do not test, a \emph{read/write
dissociation}: the direction that best \emph{reads} a behaviour off the
activations may not be the direction that best \emph{writes} it back in. This
thesis makes that hypothesis explicit and tests it, for the first time, on
safety behaviours rather than on truth.

\subsection*{Contributions}
\begin{itemize}
  \item \textbf{The first test of the read/write dissociation on safety
  behaviours}: for each behaviour the same activations yield mass-mean,
  logistic-regression and attention probes for detection and mass-mean,
  logistic-regression and contrastive-activation-addition (CAA) directions for
  steering, evaluated under one protocol, so that ``which direction detects
  best'' and ``which direction steers best'' can be compared directly.
  \item \textbf{A study across five safety behaviours}: deception and
  high-stakes awareness in this draft, and sycophancy, user intention and
  refusal in the full thesis, with the learned directions additionally compared
  against probes trained on sparse-autoencoder latents to identify which latents
  carry each behaviour.
  \item \textbf{A layer-selection method for behaviours lacking a token-level
  readout}: a normalized-indirect-effect sweep with a baseline-separability
  gate, standing in for the interchange-patching localisation used by
  Marks and Tegmark where a one-token minimal pair is unavailable.
  \item \textbf{A reusable pipeline that makes the study uniform}: each behaviour
  is a single specification (data, token masking, readout format, steering mode)
  run through identical layer selection, probe training, in- and
  out-of-distribution evaluation, and calibrated steering, so results are
  comparable across behaviours and new behaviours are added as configuration.
\end{itemize}

Our code, datasets, and an interactive user interface are available at \url{https://github.com/NimunB/probing-safety-behaviours}.

\section{Related Work}

\paragraph{Linear representations and probing}
The hypothesis that models represent features as directions in activation
space~\cite{elhage2022} motivates linear probing~\cite{alain2018}. Probes have
been trained for truthfulness~\cite{azaria2023,marks2024}, deception
\cite{goldowsky2025}, and hazardous or high-stakes content~\cite{mckenzie2025}.
Generalisation of such probes is contested~\cite{levinstein2024}, which is why
we foreground out-of-distribution evaluation throughout.

\paragraph{Mass-mean versus logistic regression}
Marks and Tegmark~\cite{marks2024} introduce mass-mean probing and show its
directions are more causally implicated than those of logistic regression or
contrast-consistent search, despite comparable classification accuracy; they
localise truth to a small group of ``summarisation'' hidden states via
interchange patching and evaluate interventions strictly out-of-distribution.
We adopt their plain (uncentred-covariance) mass-mean direction and their
mean-gap steering calibration, and we test whether their causal-superiority
finding transfers to deception and high-stakes.

\paragraph{Detecting strategic deception}
Apollo~\cite{goldowsky2025} train probes on \emph{instructed pairs}: a fact
completed under an honest versus a deceptive system prompt, taking activations
over the response tokens (excluding the final few, so the probe captures intent
rather than the completed claim) and aggregating per-token scores at
evaluation. They report AUROC and recall at a low false-positive rate, and note
a spurious correlation whereby the probe partly tracks whether a scenario is
\emph{about} deception. We reuse their instructed-pairs construction and
per-token training, and evaluate out-of-distribution on the AI-Liar
dataset~\cite{pacchiardi2024}.

\paragraph{Detecting high-stakes interactions}
The high-stakes work~\cite{mckenzie2025} trains probes on a synthetic dataset of
high- and low-stakes user interactions and evaluates on curated real-world
datasets. They compare last-token, mean, softmax and learned-attention pooling,
finding sequence-pooling essential, and report AUROC and true-positive rate at
$1\%$ false-positive rate. We adopt their dataset, pooling comparison and
metrics.

\paragraph{Contrastive activation addition}
Rimsky et al.~\cite{rimsky2024} steer behaviour by adding the mean difference
between contrastive pairs, extracted at the answer position and added at every
subsequent position. We include a CAA direction as a third steering vector,
extracted at the final prompt token so that it is not numerically identical to
the mass-mean direction on paired data.

\paragraph{Sparse autoencoders}
SAEs decompose activations into sparse, putatively interpretable
latents~\cite{cunningham2023,bricken2023}. Apollo report that probes trained on
SAE latents \emph{underperform} raw-activation probes for deception. We
therefore frame our planned SAE-latent probing as \emph{interpretation} of the
learned directions (which latents carry the signal) rather than as a competing
detector.

\section{Method}

\subsection{Model and behaviour specification}
All experiments use Llama-3.1-8B-Instruct in 8-bit precision, loaded once and
shared across behaviours; probing operates on the residual stream, which is
unaffected by weight quantisation. Each behaviour is described by a single
specification object fixing: its positive class; its token-mask policy; the
constrained yes/no readout used for causal measurement; the LLM-judge rubric;
and the steering mode, additive or (reserved for refusal) directional ablation.
Adding a behaviour is thus a matter of data plus a specification.

\subsection{Data and token masking}
\textbf{Deception (positive $=$ dishonest).} We use $512$ true statements, each
wrapped in an honest and a deceptive system prompt to form a paired training
set~\cite{goldowsky2025}. At training we take mean/last activations over the
assistant's response tokens \emph{excluding the final five tokens}; at
evaluation we score all response tokens. Out-of-distribution evaluation uses
AI-Liar~\cite{pacchiardi2024}, in which the same answer is honest or deceptive
depending on the system framing; we note this evaluation is off-policy for our
model.
\textbf{High-stakes (positive $=$ high-stakes).} We use the synthetic
\emph{models-under-pressure} training split~\cite{mckenzie2025}, pooling over
all tokens of the chat-formatted user turn, with training and evaluation using
the same mask. Out-of-distribution evaluation uses two real-world
configurations (human/assistant conversations and tool-use dialogues).
Splits are \emph{group-aware}: all tone, language and style variants sharing a
scenario identifier are kept on the same side of every split and
cross-validation fold, preventing near-duplicate leakage.

\subsection{Layer selection}
Neither probing paper selects layers by patching; patching requires a one-token
minimal pair and a token-level readout, which our behaviours lack. We therefore
use two complementary, purpose-matched diagnostics, and choose the final layers
manually from the resulting plots:
\begin{itemize}
  \item \emph{Detection layer}: group-aware five-fold cross-validated AUROC of a
  mass-mean probe on mean-pooled activations, swept over all layers.
  \item \emph{Steering block}: a per-layer \emph{normalized indirect effect}
  (NIE). Following Marks and Tegmark~\cite{marks2024}, we add the
  mean-gap-calibrated mass-mean direction at a single layer inside a few-shot
  yes/no readout and measure
  $\mathrm{NIE} = (\mathrm{PD}_{\text{steered}} - \mathrm{PD}_{\text{base}}) /
  (\mathrm{PD}^{+} - \mathrm{PD}^{-})$, where $\mathrm{PD} = P(\text{YES}) -
  P(\text{NO})$ at the final position. A gate rejects the metric where the
  baseline class gap $\mathrm{PD}^{+}-\mathrm{PD}^{-}$ is near zero, i.e.\ where
  the model cannot self-report the behaviour and the ratio is undefined.
\end{itemize}

\subsection{Probe grid}
At the detection layer we train five probes: mass-mean and logistic regression,
each with last-token and mean pooling, plus a learned single-head attention
probe (three seeds). Mass-mean uses the plain difference-of-means direction;
we omit the covariance-corrected variant, which converges to the
logistic-regression direction~\cite{marks2024} and would merely duplicate it.
The mean-pooled logistic probe is trained \emph{per token} following
Apollo~\cite{goldowsky2025} (every masked token is a labelled row); because a
fixed linear probe satisfies $\overline{w\cdot x} = w\cdot\bar{x}$, this changes
only training, not evaluation, and it makes per-token score heatmaps
interpretable. Per-token training rows are uniformly subsampled to a fixed cap
for tractability; the forward pass always sees intact prompts, so no context is
lost.

\subsection{Detection metrics}
For every (probe, dataset) pair, in- and out-of-distribution, we report AUROC
with bootstrap confidence intervals; true-positive rate at $1\%$
false-positive rate (within-dataset negatives); and accuracy at a
\emph{per-probe} decision boundary frozen on the training split: the
class-mean midpoint for mass-mean (its projection is not zero-centred), and the
model's own boundary ($\text{logit}=0$; $p=0.5$) for the trained probes.

\subsection{Causal steering}
We compare three learned directions, mass-mean, logistic-regression weights,
and CAA (extracted at the final prompt token), each calibrated so that a
coefficient of $1.0$ moves the mean negative representation onto the mean
positive along that direction~\cite{marks2024}. A separate direction is
extracted and calibrated \emph{per layer} of the steering block (a documented
departure from adding a single vector across the block, motivated by the drift
of the residual basis across depth), and added at every token position
following CAA~\cite{rimsky2024}, over a coefficient sweep
$\{-2,-1,-0.5,0,0.5,1,2\}$. We steer \emph{in-distribution first} as a positive
control, distinguishing ``the direction is not causal'' from ``the direction is
causal but does not transfer,'' and then out-of-distribution as the headline
result. Each steered generation is scored two ways: by the NIE above, and by an
LLM judge (Claude Haiku) that rates the behaviour and, separately, coherence
$1$--$10$; the coherence axis guards against interventions that ``succeed'' only
by degrading the text.

\section{Results}
\TODO{All numbers below are placeholders; the first full run is pending on the
GPU environment. Structure and metrics are final.}

\subsection{Where do the behaviours live?}
The cross-validated AUROC sweep is expected to peak in early-to-middle layers
for deception (near $28\%$ depth, following~\cite{goldowsky2025}) and
mid-network for high-stakes (following~\cite{mckenzie2025}), with the NIE sweep
identifying a contiguous causally-implicated block.
\TODO{report chosen DEC/HS detection layers and steering blocks; include the
CV-AUROC and NIE-per-layer figures, and the deception interchange-patching
heatmap.}

\subsection{Linear structure}
\TODO{PCA scatter (mean-pooled, centred) plus projection-onto-probe-direction
histograms per behaviour; comment on separability and the topic confound.}

\subsection{Detection, in- and out-of-distribution}
\begin{table}[htbp]
\caption{Detection metrics (placeholder). One block per behaviour;
columns are in-distribution test and each OOD set.}
\begin{center}
\begin{tabular}{llccc}
\toprule
Probe & Eval & AUROC & TPR@1\%FPR & Acc. \\
\midrule
MM (last)     & ID / OOD & \NUM{--} & \NUM{--} & \NUM{--} \\
MM (mean)     & ID / OOD & \NUM{--} & \NUM{--} & \NUM{--} \\
LR (last)     & ID / OOD & \NUM{--} & \NUM{--} & \NUM{--} \\
LR (mean)     & ID / OOD & \NUM{--} & \NUM{--} & \NUM{--} \\
Attention     & ID / OOD & \NUM{--} & \NUM{--} & \NUM{--} \\
\bottomrule
\end{tabular}
\label{tab:detection}
\end{center}
\end{table}
We expect sequence-pooling probes to dominate last-token probes, especially for
high-stakes~\cite{mckenzie2025}, and logistic regression to be the strongest
deception detector~\cite{goldowsky2025}.
\TODO{fill Table~\ref{tab:detection}; report ROC curves.}

\subsection{Causal steering: the read/write comparison}
\begin{table}[htbp]
\caption{Steering effect (placeholder). NIE and judged behaviour shift at
matched coefficient, per direction, in- and out-of-distribution.}
\begin{center}
\begin{tabular}{llcc}
\toprule
Direction & Setting & NIE & Judge $\Delta$behaviour \\
\midrule
Mass-mean & ID / OOD & \NUM{--} & \NUM{--} \\
Logistic  & ID / OOD & \NUM{--} & \NUM{--} \\
CAA       & ID / OOD & \NUM{--} & \NUM{--} \\
\bottomrule
\end{tabular}
\label{tab:steering}
\end{center}
\end{table}
The central test: if logistic regression tops Table~\ref{tab:detection} while
mass-mean tops Table~\ref{tab:steering}, the read/write dissociation holds for
safety behaviours as it does for truth~\cite{marks2024}.
\TODO{fill Table~\ref{tab:steering}; include dose-response curves (behaviour and
coherence versus coefficient) and qualitative steered generations.}

\section{Discussion}
\TODO{Interpret the read/write comparison; whether pooling and layer findings
replicate at 8B; the practical implication that a good monitor need not be a
good control knob, and vice versa.}

\subsection{Limitations}
The AI-Liar evaluation is off-policy for our model, so deception transfer is a
lower bound relative to on-policy evaluation~\cite{goldowsky2025}. We study a
single 8B model, below the scale of the reference works. The LLM judge is
single-source and should be validated against a small hand-labelled set. PCA
separation can reflect topic confounds rather than the behaviour itself
\cite{goldowsky2025}, and we retain no analogue of the \emph{likely} dataset
that Marks and Tegmark use to rule out a probability-versus-truth confound
\cite{marks2024}.

\section{Extending the Study}
The two extensions below are core aims of the thesis, built on the results this
draft establishes rather than left as speculative future work. The pipeline
above makes them cheap to run; the contribution is the additional evidence they
provide, not the tooling.

\paragraph{More behaviours}
Each new behaviour is a specification, not new code.
\emph{Sycophancy} fits the contrastive-pair mould directly, with an A/B
answer-choice readout replacing the yes/no readout (the readout target tokens
are already configurable). \emph{User intention}, for example benign versus
manipulative or harmful intent behind a request, uses a full-sequence mask over
the user turn, closely mirroring the high-stakes recipe. \emph{Refusal} is the
motivating case for the steering mode being abstracted: following the
refusal-direction literature, it is best manipulated by \emph{directional
ablation}, projecting the direction out of the residual stream rather than
adding to it, and the steering hook already supports this mode.

\paragraph{Probing on sparse-autoencoder latents}
Because activations are cached to disk, the SAE analysis is post-processing. We
encode the cached residual-stream activations with an SAE for this model and
train the same probe grid on its \emph{latents} rather than on the raw residual
stream, asking whether a small set of interpretable latents suffices to detect
each behaviour and which latents carry the signal. Given that probes on SAE
latents underperform raw-activation probes for deception~\cite{goldowsky2025},
the contribution here is interpretive, identifying \emph{which} latents encode
deception and high-stakes, rather than a bid to beat the raw-activation
detector.

\section{Conclusion}
\TODO{One paragraph once results exist: state whether the read/write
dissociation holds across both safety behaviours, and what that implies for
choosing between probe-based monitoring and probe-based control.}

\section*{Acknowledgment}
This work builds on the open datasets and code released by the authors of
\cite{marks2024,goldowsky2025,mckenzie2025,rimsky2024} and on the ARENA
educational materials.

\begin{thebibliography}{00}
\bibitem{marks2024} S. Marks and M. Tegmark, ``The geometry of truth: Emergent
linear structure in LLM representations of true/false datasets,'' in
\emph{Conf. on Language Modeling (COLM)}, 2024.
\bibitem{goldowsky2025} N. Goldowsky-Dill, B. Chughtai, S. Heimersheim, and
M. Hobbhahn, ``Detecting strategic deception using linear probes,''
arXiv:2502.03407, 2025.
\bibitem{mckenzie2025} A. McKenzie, U. Pawar, P. Blandfort, W. Bankes,
D. Krueger, E. S. Lubana, and D. Krasheninnikov, ``Detecting high-stakes
interactions with activation probes,'' in \emph{NeurIPS}, 2025.
\bibitem{rimsky2024} N. Rimsky, N. Gabrieli, J. Schulz, M. Tong, E. Hubinger,
and A. Turner, ``Steering Llama 2 via contrastive activation addition,'' in
\emph{Proc. ACL}, 2024.
\bibitem{alain2018} G. Alain and Y. Bengio, ``Understanding intermediate layers
using linear classifier probes,'' arXiv:1610.01644, 2018.
\bibitem{azaria2023} A. Azaria and T. Mitchell, ``The internal state of an LLM
knows when it's lying,'' arXiv:2304.13734, 2023.
\bibitem{levinstein2024} B. A. Levinstein and D. A. Herrmann, ``Still no lie
detector for language models,'' 2024.
\bibitem{elhage2022} N. Elhage et al., ``Toy models of superposition,''
\emph{Transformer Circuits Thread}, 2022.
\bibitem{cunningham2023} H. Cunningham, A. Ewart, L. Riggs, R. Huben, and
L. Sharkey, ``Sparse autoencoders find highly interpretable features in
language models,'' arXiv:2309.08600, 2023.
\bibitem{bricken2023} T. Bricken et al., ``Towards monosemanticity: Decomposing
language models with dictionary learning,'' \emph{Transformer Circuits Thread},
2023.
\bibitem{pacchiardi2024} L. Pacchiardi et al., ``How to catch an AI liar:
Lie detection in black-box LLMs by asking unrelated questions,'' in \emph{ICLR},
2024.
\end{thebibliography}

\end{document}
